\section{Stopped interfaces and signed slit data}\label{if:section}

We now turn the complete loop limit into the boundary information needed
for a conditional target law. The exploration retains its full spin
transcript, while the geometric datum records only the boundary visible
from the target. The next section will replace that boundary uniformly
over nearby transcripts. All distances in this and the next two sections
are macroscopic: the lattice is $\delta\mathbb Z^2$, with
$\delta\downarrow0$.

\subsection{Band geometry and the information revealed}

Write $Q_r=\{z\in\mathbb R^2:\|z\|_\infty<r\}$ and
$\gamma_r=\partial Q_r$.  Fix $0<u<v$ and positive clearances from the
band $\overline Q_v\setminus Q_u$ to the target and to the boundary of
the ambient chart.  There are two versions of the construction.
\begin{enumerate}
\item In the inward version, initially reveal the spins outside $Q_v$,
      and explore from $\gamma_v$ towards $\gamma_u$.  The target is a
      fixed compact subset of $Q_u$; in applications it is contained in
      a strictly smaller concentric square.
\item In the outward version, initially reveal the spins in $Q_u$,
      and explore towards $\gamma_v$.  There is a more distant fixed
      square wall, and the target lies between $\gamma_v$ and that wall,
      with positive clearances from both.  The distant wall has prescribed
      signs, including the choice of a constant minus wall.
\end{enumerate}
The statement also allows a compact interval of choices for each radius,
provided the intervals are separated and all the stated clearances hold
uniformly. For a requested radius $r$, choose deterministically a nearest
$r_\delta\in\delta\mathbb N$, with ties resolved by a fixed rule, and
use the contour $\partial Q_{r_\delta}$. For sufficiently small $\delta$
these rounded radii preserve all the prescribed separations, and their
contours follow primal grid edges. The
interpolated zero set meets such a contour only at edge midpoints:
it never passes through a spin vertex, and the adjacent triangles
lie on opposite sides of the contour edge.  Thus these meetings are
crossings, including near the square corners.
In either version include the vertices of the rounded starting contour
in the initial revealed layer.
For measurability statements, the \emph{band spins} include the vertices
of triangles meeting the closed rounded band. This convention adds only
an $O(\delta)$ layer and includes the starting-contour vertices.

\begin{remark}[Thin fixed collars]\label{if:thin-collars}
No universal lower bound on $v/u-1$ is required. Each collar has
strictly positive width fixed before the mesh limit, and the constants
and mesh threshold may depend on that width. For example, collinear
segments of length $\ell$, separated by gaps $g>0$, fit in disjoint
squares with containing halfside $\ell/2+g/8$ and farther-wall
halfside $\ell/2+g/4$. There is room between these squares for any
fixed finite list of nested contour levels with positive clearances.
When pieces come from a polygonal contour, first remove neighborhoods
of corners and boundary intersections, then take segment lengths
smaller than those separations. The remaining straight pieces admit
the same construction. Its constants can depend on $g/\ell$; a
fixed enlargement factor independent of $g/\ell$ is unnecessary.
\end{remark}

When radii vary, reserve one fixed larger permitted window containing
their compact ranges, and choose all common source and test contours
with positive clearance from the entire range of possible stopping
bands.  The regular barriers constructed below stay on the fresh side
of the realized visible wall, with the prescribed target clearance;
they need not lie between each realization's numerical radii $u,v$.
The reserved window provides room for their offsets and tip caps.

We encode lattice boundary data on the requested continuum contours.
Choose disjoint fixed intervals about the two permitted radius ranges,
with positive endpoint clearances. Let $f_\delta:[0,\infty)\to[0,\infty)$
be the increasing piecewise-linear map that sends $u_\delta$ to $u$ and
$v_\delta$ to $v$, fixes the interval endpoints, and is the identity
outside those intervals, linear between consecutive specified knots.
Define
\[
 R_\delta(0)=0,\qquad
 R_\delta(z)=\frac{f_\delta(\|z\|_\infty)}{\|z\|_\infty}z
 \quad(z\ne0).
\]
This homeomorphism and its inverse move every point by $O(\delta)$,
uniformly over the permitted radii: on each affine interval the error
$f_\delta(r)-r$ interpolates endpoint errors of size at most
$\delta/2$. It maps the closed rounded band onto the exact band and
each rounded contour onto its requested counterpart. We apply
$R_\delta$ only when recording the geometry of the exploration below.
The physical spins, reveal, and Gibbs laws remain on the original lattice.

For the probability limit, sample spins in a fixed larger square with
plus boundary conditions.  Other buffered laws will enter only through
Corollary~\ref{fd:buffered-ratio}.  The sign-reversed reference square may
equally well be used.  A finite collection of band observations is always
computed from one common spin configuration, so that its joint law has
an unambiguous meaning.

Fix one diagonal in every square face, without adding interactions to
the Ising Hamiltonian.  Affinely interpolate the spins in each triangle.
The zero set consists of noncrossing polygonal strands.  After the initial
reveal, inspect every strand meeting the starting contour.  Follow it on
the fresh side until its first return to the starting contour or its
first hit of the other contour.  In the latter case also follow the
strand from its other starting endpoint, when there is one, stopping by
the same rule.  Every time a triangle is traversed, inspect its vertex
spins.  The order of exploring the starting endpoints is fixed
deterministically.  Repeated inspections reveal no additional information.

From the target, the visible boundary consists of a base formed by the
starting contour and returning excursions, together with strands stopped
at the second contour and their two side signs. This is the
target-facing datum formalized in Definition~\ref{if:datum}.

\begin{lemma}[Stopping property]\label{if:stopping}
For each $\delta>0$, this is a finite adaptive spin reveal.  Conditional
on its complete information, the unexamined spins have the Ising Gibbs
law with the examined spins fixed.  The revealed set lies on the starting
side and in an $O(\delta)$ enlargement of the band.  The traced strands,
their side signs, and the target-facing boundary datum defined below
can be computed from the band spins alone.
\end{lemma}

\begin{proof}
There are finitely many starting edges and finitely many triangles.
At a step, the next triangle is determined by the previously exposed
zero segment and the fixed diagonal convention.  Its spins determine
the continuation.  Hence the event that a specified reveal transcript
occurs depends only on the spins recorded in that transcript.  In the
finite Gibbs density, conditioning on such a transcript fixes precisely
those spins and leaves the original interaction weights on all remaining
variables.  This proves the stopping assertion directly, including the
case of a random examined set.

The trace stops at the second contour and examines only adjacent
triangles, which proves the location assertion.  Although the initial
reveal contains an entire side of the band, the algorithm deciding a
strand within the band uses only the spins adjacent to the starting
contour and in the band itself.  The target-facing datum is determined
by these traces and their local signs.  No spins in the rest of the
initially revealed side are needed to compute it.
\end{proof}

\subsection{The visible Jordan base and its slits}

The following definitions refer first to an oriented collection of
simple disjoint continuum loops.  Restrict a loop to the fresh side of
the starting contour.  A maximal such portion returning to that contour
before reaching the second contour is a \emph{returning chord}.  A
portion reaching the second contour is retained only up to its first
hit there and is a \emph{stopped strand}.  Cut the fresh region along
these curves and take the component containing the untraced region
beyond the second contour.  Retain only boundary pieces seen from this
component.  The same definition is made for the resolved lattice walls.

The boundary is read from the remaining domain: a traversal follows
one side of each slit to its tip and returns along the other.
Thus the two sides of a slit are distinct boundary arcs even though
they have the same geometric image. This is the prime-boundary
traversal used below.

\begin{definition}[Admissible facing datum]\label{if:datum}
A facing datum consists of an oriented Jordan base, finitely many
pairwise disjoint simple slits attached to it at distinct roots, the
root and tip marks, and signs on the two sides of each slit and on
the intervening base arcs.  Signs on a slit are opposite.  In the
prime-boundary traversal the signs change at its tip and continue
unchanged through each occurrence of its root.  The whole datum lies
in the closed band, and its tips lie on the second contour.  If there
are no slits, the base has one constant side sign.

The lattice datum is the image of its physical boundary record under
$R_\delta$, retaining roots, tips, orientation, and side signs.
When radii vary over permitted windows,
$\mathcal X_{\mathrm{band}}$ denotes the union of admissible data from
those radius choices. The numerical radii are not extra marks; for
example, a datum without slits need not determine its second radius.

Within a fixed finite combinatorial type, use the maximum uniform
distance modulo increasing parametrizations of the base and slits,
respecting root attachments, side labels, and orientation.  Cap this
distance at one and put distance one between different types.  Denote
this metric space by $\mathcal X_{\mathrm{band}}$.

For approximation centers we also use an ambient space
$\widehat{\mathcal X}$ of finite marked continuous walks of these
combinatorial types. Retain the orientations, side labels, root and tip
marks, and shared coordinates of attached roots, but drop simplicity,
disjointness, and the exact band-containment and contour-incidence
requirements. Use the same matching-distance formula, identifying
representatives at distance zero. Then
$\mathcal X_{\mathrm{band}}\subset\widehat{\mathcal X}$ has its
stated metric as the induced metric. A ball with an ambient prototype
as center always means its intersection with
$\mathcal X_{\mathrm{band}}$ when used as a set of facing data.
\end{definition}

\begin{lemma}[Visible data and their continuity]\label{if:visible-data}
For almost every choice of the two contour radii, the continuum
construction from nested $\mathrm{CLE}_3$ yields an admissible facing
datum.  Its base may contain infinitely many returning chords but is
a Jordan curve.  The lattice data converge to it in
$\mathcal X_{\mathrm{band}}$.  These assertions hold jointly for any
fixed finite collection of bands.
\end{lemma}

\begin{proof}
We give the deterministic continuity argument on the probability-one
event of simple, disjoint loops and local finiteness at every positive
diameter.  For a continuous real function on a compact parameter
interval, its local maximum and minimum values form a countable set:
each is the maximum or minimum on an interval with rational endpoints
contained in a witnessing neighborhood.  Apply this to the square
radius along each loop.  There are countably many loops, so the union
of exceptional levels is countable for each collection.  Fubini's
theorem permits a deterministic full-measure set of radius choices for
the random collection.  At the selected levels there is no one-sided
contact or constant interval.

Only finitely many excursions have diameter greater than a prescribed
$a>0$.  Indeed, only finitely many loops have that diameter; on each,
uniform continuity of a parametrization gives a positive minimum
parameter length for such an excursion.  Their parameter intervals
are disjoint.  A stopped strand has diameter at least the positive
distance between the contours, so there are finitely many stopped
strands.  Their tips are distinct.  The two stopped ends of the same
deep excursion could have the same tip only if that tip were a
one-sided contact at the second level, which was excluded.

Distinct excursion endpoints on the starting contour are also
distinct.  Endpoints on different loops cannot coincide by
disjointness.  A shared endpoint for two excursions of one simple
loop would give a one-sided level contact.  Excursion interiors miss
the starting contour, and simplicity and disjointness exclude all
other intersections.

Each visible returning chord deletes the open interval of the starting
contour which is cut off from the target component.  Visible deletion
intervals have disjoint interiors: in a pair of nested deletions the
inner chord is screened and is not visible.  A small chord deletes a
small interval.  To see this, join its endpoints by the short arc of
the starting square.  As the chord diameter tends to zero, uniform
continuity of the inverse parametrization of that square makes this
arc small.  The resulting small Jordan pocket cannot contain the
target beyond the second contour, which is separated by a fixed
positive distance.  The erased interval must therefore be this short
arc, not the complementary long arc.  This argument works also on
the exterior side; equivalently, compactify the exterior by a fixed
inversion away from the band.

Parametrize the starting contour and replace every deleted interval
by a parametrization of its chord with the same endpoints.  Retain
the original parametrization elsewhere.  At an accumulation of
replacement intervals, the interval diameters and chord diameters
tend to zero, by the preceding finiteness and small-pocket facts.
The resulting map is consequently continuous.  It is injective:
chord interiors miss the starting contour and one another, their
endpoints are distinct, and retained points of the original contour
remain distinct.  A continuous injection of the compact circle into
the plane is a Jordan parametrization.  This proves, in particular,
that infinitely many small returning chords do not create an
accumulating pinch.  Adding the finitely many visible stopped strands
gives the asserted base and slits.

For continuity, first choose $a>0$ outside the countable set of excursion
diameters and retain the finitely many excursions of diameter greater
than $a$. Uniform convergence of the parent loops
gives convergence of these excursions, their endpoints, and their
first hits of the second contour.  Here is the role of the generic
levels.  A compact subarc strictly on one side of a level stays there
under a small uniform perturbation. Arbitrarily near an excursion
endpoint on its other side in parameter time, the limiting loop has
points on the opposite side of the level; otherwise the endpoint
would be a one-sided contact. These points bracket the perturbed endpoint
between arbitrarily nearby parameter values.  The same argument at
the cutoff brackets its first hit.  Splitting, merging, or losing a
retained excursion would contradict one of these bracketing
properties.

The circular order of the finitely many retained endpoints is stable.
Their visibility is stable as well.  An additional chord of small
diameter cannot screen a retained chord: its deleted pocket has small
diameter and cannot contain that retained chord.  After matching the
retained pieces, match the remaining intervals of the bases in their
circular order.  All discrepancies on these intervals have diameter
tending to zero as $a\downarrow0$.  This constructs parametrizations
with uniformly small distance.  Include the finitely many slit roots
in the matching; their stopped strands already converge by the same
bracketing argument. In particular the fixed band width and the generic
second level give eventual equality of the finite stopped-strand list,
with its cyclic root and tip order. Only the returning chords require
the subsequent limit $a\downarrow0$. The recording map $R_\delta$ adds
an error of $O(\delta)$ and ensures that the lattice records have the
exact contour incidences required by the data space.

At the lattice level, a returning chord replaces an interval whose
two surviving side signs agree.  The only sign changes in a traversal
of the visible prime boundary occur at stopped tips.  Thus microscopic
returning chords introduce no additional alternations between tips.
When slits are present, their oriented sides determine the signs on
all intervening base arcs, so these signs converge by
Lemma~\ref{if:resolution-convergence}.

To recover the constant sign when there are no tips, we use one further
property of nested $\mathrm{CLE}_3$: a deterministic interior point is
almost surely surrounded by infinitely many nested loops
\cite[Section~1, before Proposition~1]{SSW2009}. Their diameters tend to
zero by local finiteness. Choose such a point on the starting contour.
A sufficiently small surrounding loop misses the second contour and
has a returning excursion into the fresh side. If this chord is hidden,
any chord screening it has diameter bounded away from zero by the
small-pocket argument. There are only finitely many such screening
chords, so one is visible. Its oriented parent loop supplies the constant
side sign. That loop and its finitely many ancestors converge with their
orientations by Lemma~\ref{if:resolution-convergence}. This proves sign
convergence also in the no-slit case. The additional fixed-point nesting
property is used only for this color anchor, not for the preceding
deterministic construction of the Jordan base.

The deterministic argument applies simultaneously to any fixed finite
list.  Lemma~\ref{if:resolution-convergence} and the continuous mapping
theorem finish the proof.
\end{proof}

\begin{remark}\label{if:metric-separable}
The ambient space $\widehat{\mathcal X}$ in Definition~\ref{if:datum}
is separable. Approximate finitely many parametrized walks by rational
piecewise-linear walks, approximating shared endpoints and roots
together. These prototypes need not be simple or satisfy the exact
band constraints; they form a countable dense set in the ambient
matching metric. The countable union over finite combinatorial types
remains separable. An ambient countable base restricts to a countable
base on $\mathcal X_{\mathrm{band}}$, so that subspace is also
separable. Only realizable admissible data, not the prototypes, are
used to construct Gibbs boundaries.
\end{remark}
